% Standalone problem document, generated from the adjacent README.md.
% Compile with XeLaTeX. Mathematical content is maintained in Markdown.
\documentclass[11pt,a4paper]{article}
\usepackage[fontsize=10.5pt]{scrextend}
\usepackage[margin=22mm,headheight=15pt,headsep=9mm,footskip=11mm]{geometry}
\usepackage{fontspec}
\setmainfont{texgyrepagella}[Extension=.otf,UprightFont=*-regular,BoldFont=*-bold,ItalicFont=*-italic,BoldItalicFont=*-bolditalic]
\setsansfont{texgyreheros}[Extension=.otf,UprightFont=*-regular,BoldFont=*-bold,ItalicFont=*-italic,BoldItalicFont=*-bolditalic]
\setmonofont{texgyrecursor}[Extension=.otf,UprightFont=*-regular,BoldFont=*-bold,ItalicFont=*-italic,BoldItalicFont=*-bolditalic,Scale=MatchLowercase]
\usepackage{amsmath,amssymb}
\usepackage{unicode-math}
\setmathfont{texgyrepagella-math.otf}
\usepackage{xcolor,microtype,enumitem,needspace,fancyhdr}
\usepackage{bookmark}
\definecolor{ink}{HTML}{17344B}
\definecolor{accent}{HTML}{197D80}
\definecolor{muted}{HTML}{596773}
\hypersetup{colorlinks=true,linkcolor=accent,urlcolor=accent,pdftitle={RA-17 - Minimum
linear measurements for uniform recovery of real low-rank
matrices},pdfauthor={Open Problems in Numerical Linear Algebra}}
\urlstyle{same}
\setlength{\parindent}{0pt}
\setlength{\parskip}{5pt plus 1pt minus 1pt}
\linespread{1.04}
\setlength{\emergencystretch}{3em}
\widowpenalty=10000
\clubpenalty=10000
\displaywidowpenalty=10000
\allowdisplaybreaks[1]
\setlist{nosep,leftmargin=1.5em}
\providecommand{\tightlist}{\setlength{\itemsep}{0pt}\setlength{\parskip}{0pt}}
\providecommand{\pandocbounded}[1]{#1}
\pagestyle{fancy}
\fancyhf{}
\fancyhead[L]{\sffamily\footnotesize\color{muted}OPEN PROBLEMS IN NUMERICAL LINEAR ALGEBRA}
\fancyhead[R]{\sffamily\bfseries\footnotesize\color{accent}RA-17}
\fancyfoot[L]{\parbox[b]{0.9\textwidth}{\sffamily\fontsize{8}{10}\selectfont\color{muted}Editorial ratings; a bounded search cannot certify that a problem remains unsolved.\\Literature check: 2026-09-13}}
\fancyfoot[R]{\sffamily\footnotesize\color{muted}\thepage}
\renewcommand{\headrulewidth}{0.3pt}
\renewcommand{\headrule}{\hbox to\headwidth{\color{accent}\leaders\hrule height \headrulewidth\hfill}}
\makeatletter
\renewcommand{\section}{\Needspace{5\baselineskip}\@startsection{section}{1}{0pt}{12pt plus 2pt minus 2pt}{4pt}{\sffamily\bfseries\large\color{ink}}}
\renewcommand{\subsection}{\Needspace{4\baselineskip}\@startsection{subsection}{2}{0pt}{9pt plus 1pt minus 1pt}{3pt}{\sffamily\bfseries\normalsize\color{ink}}}
\makeatother
\setcounter{secnumdepth}{0}
\begin{document}
{\sffamily\small\color{muted}Randomized and low-rank approximation\par}
\vspace{3pt}
{\raggedright\hyphenpenalty=10000\exhyphenpenalty=10000\sffamily\bfseries\fontsize{21}{24}\selectfont\color{ink}Minimum
linear measurements for uniform recovery of real low-rank matrices\par}
\vspace{7pt}
{\sffamily\small\textbf{Difficulty:} challenging\quad\textbf{Importance:} interesting
to the community\par}
{\sffamily\small\bfseries\color{accent}Status: Partially resolved\par}
\vspace{4pt}
{\color{accent}\hrule height 0.8pt}
\vspace{6pt}
\textbf{Rating rationale:} The all-dimension real measurement
classification needs new control of exceptional kernels and topological
obstructions, hence challenging. Its main impact is on the community
working on low-rank matrix sensing and recovery.

\subsection{Problem statement}\label{problem-statement}

For integers \(d\geq2\) and \(1\leq r\leq\lfloor d/2\rfloor\), let

\[\mathcal M_{d,r}(\mathbb R)=\{X\in\mathbb R^{d\times d}:\mathop{\mathrm{rank}}\nolimits X\leq r\}.\]

For real measurement matrices \(A_1,\ldots,A_m\), define

\[\mathcal A(X)=
(\mathop{\mathrm{tr}}\nolimits(A_1^\top X),\ldots,\mathop{\mathrm{tr}}\nolimits(A_m^\top X)).\]

Determine the exact integer

\[\mu_{\mathbb R}(d,r)=
\min\{m:\exists A_1,\ldots,A_m\in\mathbb R^{d\times d},\
\mathcal A|_{\mathcal M_{d,r}(\mathbb R)}\text{ is injective}\}\]

for all such \(d,r\).

Injectivity is uniform: every two matrices of rank at most \(r\) with
equal measurements must coincide. Equivalently, \(\ker\mathcal A\) must
contain no nonzero matrix of rank at most \(2r\). No randomness,
stability guarantee, computationally efficient decoder, positive
semidefiniteness, or generic-signal exception is required. The problem
concerns unrestricted real square matrices and unrestricted real linear
measurements.

\subsection{Partial continuation --- 13 September
2026}\label{partial-continuation-13-september-2026}

\textbf{Author:} Sidney Holden, Center for Computational Biology,
Flatiron Institute, Simons Foundation
(\href{https://github.com/ajt60gaibb/OpenProblemsInNLA/blob/main/references/holden-ra17-continuation-2026-09-13/SUBMISSION.md}{verified
affiliation}).

The
\href{https://github.com/ajt60gaibb/OpenProblemsInNLA/blob/main/references/holden-ra17-continuation-2026-09-13/writeup/RA17_topological_relaxation.pdf}{continuation
manuscript}, Sections 3--4, characterizes two topological relaxations
one measurement below the complex generic count: existence of a
continuous odd nonvanishing map on the low-rank unit link, and existence
of the corresponding number of independent continuous evaluation-bundle
sections. Each is equivalent to even complex determinantal degree. These
do not establish a system of constant linear measurement matrices.
Section 5 rederives the rank-one index lower bound; Section 2 supplies
integer upper-bound measurements.

The
\href{https://github.com/ajt60gaibb/OpenProblemsInNLA/blob/main/references/holden-ra17-continuation-2026-09-13/independent-review.md}{independent
informal AI-agent review} records the checked scope and arithmetic
verification.
\href{https://github.com/ajt60gaibb/OpenProblemsInNLA/blob/main/references/holden-ra17-continuation-2026-09-13/README.md}{Source,
code and provenance}. No Lean verification or external human peer review
is asserted. This extends the earlier partial submission in
\href{https://github.com/ajt60gaibb/OpenProblemsInNLA/pull/197}{PR
\#197}; no priority claim is made.

\textbf{Remaining target:} the exact count for every allowed pair. In
particular, the submission leaves \(19\leq\mu_{\mathbb R}(6,1)\leq20\)
and \(35\leq\mu_{\mathbb R}(10,1)\leq36\) undecided. Section 7
identifies the missing existence or impossibility of a 17-dimensional
real space of six-by-six matrices with every nonzero matrix of rank at
least three. RA-17 remains \textbf{Partially resolved} and counted as
open.

\subsection{Why it matters}\label{why-it-matters}

This is the information-theoretic limit of noiseless real matrix
sensing. It distinguishes intrinsic identifiability from the larger
measurement budgets that particular algorithms or stability guarantees
may require.

\subsection{References and status
check}\label{references-and-status-check}

\begin{enumerate}
\def\labelenumi{\arabic{enumi}.}
\tightlist
\item
  Z. Xu, \emph{Signal Recovery on Algebraic Varieties Using Linear
  Samples}, \href{https://arxiv.org/pdf/2506.17572}{arXiv:2506.17572},
  §5, Problem 5.1 and Remark 5.6, pp.13--15; published in
  \href{https://doi.org/10.1007/s10114-026-5299-y}{Acta Mathematica
  Sinica \textbf{42} (2026), 740--754}.
\item
  Z. Xu, \emph{The minimal measurement number for low-rank matrix
  recovery}, \href{https://doi.org/10.1016/j.acha.2017.01.005}{Applied
  and Computational Harmonic Analysis}, the original real counterexample
  and complex exact-count results cited as reference 18 in the survey.
\end{enumerate}

The 2026 survey explicitly retains the real problem. It gives
\(\mu_{\mathbb R}(d,r)\leq4dr-4r^2\) and equality in specified dimension
families, while its real \((d,r)=(4,1)\) construction uses 11
measurements rather than the complex threshold 12. Thus the uniformly
asserted equality with the complex count is already false and is not
proposed here. The complex count itself is solved. Current version
checks and searches for Xu, the exact title, and 2025--2026 minimal real
low-rank measurements found no general resolution.

\subsection{Audit --- 2026-09-10}\label{audit-2026-09-10}

Independently rechecked \href{https://arxiv.org/html/2506.17572v2}{Xu,
Theorem 5.5 and Remark 5.6}, the unchanged June 2025 arXiv v2 record,
and the
\href{https://actamath.cjoe.ac.cn/Jwk_sxxb_en/EN/10.1007/s10114-026-5299-y}{2026
publication}. Equality is proved when \(d=2^k+r\) (subject to the stated
rank range) or \(d=2r+1\), so Partially resolved is warranted. Targeted
later searches found no full real measurement-count classification.
Changed extreme/broad ratings to challenging/community, matching the
concrete exceptional-kernel obstacle and the scope of noiseless matrix
sensing. The unresolved aim is a general sharp formula or classification
in \(d,r\); individual fixed cases can already be expressed as decidable
real polynomial feasibility problems.
\end{document}
